\newpage

\section{Conclusion}
\label{ch:conclusion}

The expansion of structured knowledge graphs has made vast amounts of highly curated data available to the public and enterprise domains. However, accessing this information has remained a persistent challenge due to the technical complexity of the SPARQL query language and the deep semantic mismatches between natural language and rigid graph ontologies. Traditional Text-to-SPARQL approaches—relying on supervised fine-tuning, sequence-to-sequence translation, or rule-based templates—have demonstrated significant limitations. Their inherent brittleness, lack of transparency, and dependency on domain-specific training prevent them from scaling dynamically to massive, evolving environments like Wikidata.

This thesis set out to bridge the gap between natural language and the Semantic Web by investigating and implementing a novel, LLM-driven agentic framework. By moving beyond linear pipelines, this research introduced a stateful, autonomous workflow capable of iterative planning, execution, and semantic self-correction.

\subsection{Summary of Contributions}
To address the research gaps identified in the literature, this work successfully designed and deployed a multi-stage architecture utilizing LangGraph. The core contributions of this methodology include:
\begin{enumerate}
    \item \textbf{Hybrid Semantic RAG:} A dynamic schema-linking strategy that combines lexical search with dense vector embeddings (via Qdrant), allowing the system to operate in a completely zero-shot, domain-agnostic capacity without requiring prior task-specific pretraining.
    \item \textbf{Topological Grounding:} A runtime graph discovery mechanism that retrieves 1st-hop incoming and outgoing edges for candidate entities, grounding the LLM's reasoning in the actual structural reality of the endpoint to prevent schema hallucination.
    \item \textbf{Execution-Guided Refinement:} An autonomous validation loop that utilizes actual endpoint feedback (syntax errors, empty result sets) to diagnose failures, adjust semantic constraints, and iteratively correct generated queries in real-time.
\end{enumerate}

\subsection{Key Findings}
Empirical evaluation on the highly complex, multilingual QALD-10 dataset (via the GERBIL QA platform) demonstrated the robustness and versatility of the proposed solution. The agentic workflow achieved a peak Macro F1 score of 0.50 on the English dataset using the \texttt{gpt-4.1-mini} backbone, proving that dynamic RAG and autonomous refinement can effectively substitute for rigid supervised learning.

Furthermore, comparative analysis revealed that efficiently scaled open-weight models, such as \texttt{nemotron-3-nano-30b}, exhibit surprising linguistic robustness when integrated into a tool-augmented loop, outperforming proprietary models in specific European languages (German and Russian). 

Most significantly, the ablation study isolated the impact of topological grounding, revealing a \textbf{30\% relative increase} in Entity Linking precision (reaching up to 82\%). This finding fundamentally redefines the current bottleneck in Knowledge Graph Question Answering (KGQA): the primary challenge is no longer identifying \textit{which} entity is mentioned, but rather overcoming the \textbf{Schema Gap}—the complex structural relationships and multi-hop schema mappings required to connect those entities syntactically in SPARQL.

\subsection{Future Work}
While the agentic framework establishes a highly interpretable and scalable foundation for semantic parsing, several avenues remain for future research to close the performance gap between high entity linking precision and final query execution accuracy.

Future work will focus on:
\begin{itemize}
    \item \textbf{Scaling the Semantic Index:} Expanding the localized vector retrieval index to support continuous, automated integration with the massive, ever-evolving Wikidata ontology, ensuring seamless zero-shot adaptation to newly added entities and properties.
    \item \textbf{Advanced Subgraph Injection:} Developing intelligent truncation and multi-hop graph summarization techniques to handle "hub" entities, ensuring that the LLM's context window is not saturated by excessively connected nodes while still preserving critical relational pathways.
    \item \textbf{Re-ranking and Answer Selection:} Augmenting the Validation Node with advanced, state-of-the-art answer-selection mechanisms to better evaluate edge-case queries that return logically valid but semantically misaligned result sets.
\end{itemize}

Ultimately, by transforming ambiguous natural language questions into semantically rich, verifiable SPARQL queries, this research advances the interaction between natural language processing and Semantic Web technologies. The integration of LLM agents with topological discovery paves the way for more intelligent, user-friendly, and adaptable tools for navigating the world's most complex knowledge systems.
